% !TeX program = xelatex
\documentclass[11pt,a4paper]{ctexart}
\usepackage[margin=25mm,headheight=15pt]{geometry}
\usepackage{amsmath,amssymb,amsthm,mathtools,booktabs,longtable,array,fancyhdr,xcolor,pdflscape}
\usepackage[unicode,colorlinks=true,linkcolor=blue!45!black,citecolor=blue!45!black,urlcolor=blue!45!black]{hyperref}
\newtheoremstyle{cnplain}{7pt}{7pt}{\normalfont}{}{\bfseries}{。}{0.5em}{}
\theoremstyle{cnplain}\newtheorem{theorem}{定理}[section]
\newtheorem{proposition}[theorem]{命题}
\newtheorem{lemma}[theorem]{引理}
\newtheorem{corollary}[theorem]{推论}
\theoremstyle{definition}\newtheorem{definition}[theorem]{定义}
\theoremstyle{remark}\newtheorem{remark}[theorem]{注}
\renewcommand{\proofname}{证明}
\newcommand{\Q}{\mathbb Q}\newcommand{\cC}{\mathcal C}
\newcommand{\cS}{\mathcal S}\newcommand{\dR}{\mathrm{dR}}
\newcommand{\rk}{\operatorname{rk}}\newcommand{\im}{\operatorname{im}}
\setlength{\emergencystretch}{3em}\allowdisplaybreaks[2]
\pagestyle{fancy}\fancyhf{}
\fancyhead[L]{\small 无质量费曼图：MZV与物理读出}
\fancyhead[R]{\small 研究主线}\fancyfoot[C]{\thepage}
\title{从图结构预报MZV：无质量串并联族\\[5pt]
\large 全阶维数展开、局部减除与banana的正谱}
\author{费曼周期研究项目\quad 物理分类证明补篇}
\date{2026年10月14日}
\begin{document}\maketitle
\begin{abstract}
给定无质量双端串并联分解，逐层一圈卷积在不求完整多圈积分的前提下
保证：规定Euler规范后的全部维数展开系数属于单个zeta值生成的代数。
无质量banana提供完整物理实例，反项、MS有限部分、MOM函数与正谱均有全边数公式。
四线谱曲面的有理化说明同一拓扑怎样在全零质量层脱离正质量K3。
本文包含七项带证明陈述，区分常数MZV、尺度对数、函数表示与原生动机。
经典积分恒等式不主张新颖性；一般森林结论保留明确闭合与IR条件。
\end{abstract}
\tableofcontents\clearpage
\section{无质量图的MZV肯定规则：串并联、减除与同拓扑退化}
\label{sec:masslessmzv}

本节补分类表的肯定方向：在完整积分以前，从图的分解证书预报所需常数。
接受域是无质量、非零欧氏外动量的标量两点图；物理对象是维数正规化的
Laurent系数、在明确条件下的局部反项，以及无质量banana的完整重整化函数。
这里不假定每个图都属于同一个可重整化拉格朗日量；图权、耦合和外腿归一化另行乘回。
一圈卷积和单尺度子图替换是经典规则\cite[式(1.2.1)、\S5.2.1]{PanzerMZV}。
以下自给推导它们对分类、物理减除及特殊质量层的具体结论，不主张经典公式的新颖性。

\subsection{物理输入与一圈卷积}
令$D=4-2\varepsilon$，$q^2=x>0$，每圈测度为$d^Dk/\pi^{D/2}$。
边的传播子为$(k_e^2)^{-a_e}$，初始$a_e\in\mathbb Z_{>0}$。
所有外动量只在两个指定端点注入；其余顶点总外动量为零。
若图有$h$圈，使用明确的Euler归一化
\[
 \widehat F_G(x,\varepsilon)
 =e^{h\gamma_E\varepsilon}\mu^{2h\varepsilon}F_G(x,\varepsilon),
 \qquad L=\log(x/\mu^2),\qquad
 \mathcal Z_1=\mathbb Q[\zeta(2),\zeta(3),\ldots]\subseteq\mathcal Z.
\]
此处$\mu>0$是同一减除尺度。更换测度会改变显式$\pi$与$4\pi$因子；
没有先去掉这些规范就不能声称系数属于普通MZV代数。
维数不为整数的整个积分不是一个已构造的有限维Nori对象；本节的算术结论逐个展开系数陈述。

\begin{lemma}[实际卷积核]\label{lem:masslesskernel}
在$0<\operatorname{Re} A,\operatorname{Re} B<D/2$及
$\operatorname{Re}(A+B)>D/2$的实正整数维数收敛域，
\begin{align}
 \int\frac{d^Dk}{\pi^{D/2}(k^2)^A((q-k)^2)^B}
 &=x^{D/2-A-B}\,g_D(A,B),\\
 g_D(A,B)&=
 \frac{\Gamma(D/2-A)\Gamma(D/2-B)\Gamma(A+B-D/2)}
 {\Gamma(A)\Gamma(B)\Gamma(D-A-B)}.\label{eq:masslesskernel}
\end{align}
两边以Schwinger参数形式亚纯延拓到维数和指数的一般复值。
\end{lemma}
\begin{proof}
对两条边引入正参数$t,u$，完成平方后动量高斯积分给
$(t+u)^{-D/2}\exp[-tu x/(t+u)]$。
置$t=\lambda v,u=\lambda(1-v)$，$0<v<1$。
先积$\lambda$得到$\Gamma(A+B-D/2)x^{D/2-A-B}$，
再积$v$得到$B(D/2-B,D/2-A)$，连同原来的
$1/(\Gamma(A)\Gamma(B))$即为所示公式。
收敛域中Tonelli或绝对Fubini合法；复维数版本由该参数公式定义并唯一亚纯延拓，
不是在非整数维度上假想一个通常的Lebesgue空间。
\end{proof}

\subsection{边表证书与全阶常数类}
双端串并联证书是一棵二叉树。叶是有标号边；每个结点的两个子图除所声明端点外不相交。
串联$S$仅把前者末端与后者首端粘合；并联$P$把两对端点分别粘合。
最终叶集须恰为输入边集，端点粘合须恢复原边表，并逐顶点核对无内部动量注入。
这里的并联包含任意边数banana，串联中允许二价顶点。
这是一项充分的接受条件，不是所有MZV图的必要条件。

\begin{theorem}[串并联证书预报全阶单$\zeta$代数]\label{thm:masslesssp}
给定上述证书，设$E$为边数，$h$为圈数，$A_0=\sum_e a_e-2h$。
不需要展开任一多圈积分即可断定
\[
 \widehat F_G(x,\varepsilon)
 =x^{-A_0}e^{-h\varepsilon L}\widehat C_G(\varepsilon),
 \qquad \widehat C_G\in\mathcal Z_1((\varepsilon)).
 \label{eq:masslessspclass}
\]
更精确地，证书给有限有理函数$R_G\in\mathbb Q(\varepsilon)$和有理数$c_k$使
\[
 \widehat C_G(\varepsilon)=R_G(\varepsilon)
 \exp\left(\sum_{k\ge2}\frac{\zeta(k)}k c_k\varepsilon^k\right).
 \label{eq:masslessgammaexp}
\]
若$R_G$在零点的极点阶为$p$，则第$\varepsilon^j$系数可表达为
权重至多$j+p$的单$\zeta$乘积；计$L$的权重为1，同一界适用于去掉$x^{-A_0}$后的系数。
这里是表示的权重上界，不假定$\zeta$值间的代数独立性。
\end{theorem}
\begin{proof}
每个子图记$F=Cx^{-A}$，从叶$(h,A,C)=(0,a_e,1)$开始。
串联时动量相同、环变量互相独立，因此
\[
 h=h_1+h_2,\quad A=A_1+A_2,\quad C=C_1C_2.
\]
并联时分流动量为$k,q-k$，额外增加一次真实卷积，因此
\[
 h=h_1+h_2+1,\quad A=A_1+A_2-D/2,\quad
 C=C_1C_2g_D(A_1,A_2).
\]
从而每个子图均有$A=a+h\varepsilon$、$a\in\mathbb Z$。
恒等式先在各叶指数独立的收敛域证明：指定任一子树的有效指数
$0<A<D/2$，串联可把它分成两个正数；并联可选择两个小于$D/2$的正数，
其和为$A+D/2$。反向归纳给所有叶的非空同时收敛域。
再亚纯延拓至原正整数指数和$D=4-2\varepsilon$，无需非法交换发散积分。

用$\Gamma(z+1)=z\Gamma(z)$把所有Gamma因子移至$\Gamma(1+c\varepsilon)$。
整数位移只产生有理函数。非正整数的分子Gamma具有非零斜率：
并联结点的三条分子斜率为$-(h_1+1),-(h_2+1),h_1+h_2+1$；
所以不会遇到沿整条$\varepsilon$直线恒为极点的分子。
斜率为零的分母Gamma仅来自无圈正指数子树，也不形成此问题。
各并联结点的分子斜率和减去分母斜率和恰为1。
应用经典展开\cite[5.7.3]{DLMFGamma}
\[
 \log\Gamma(1+z)=-\gamma_E z+
 \sum_{k\ge2}\frac{(-1)^k\zeta(k)}k z^k,
\]
恰有$h$个并联结点，故Euler常数项为$-h\gamma_E\varepsilon$，被规定的规范消去。
剩下$c_k=\sum_\nu\sigma_\nu(-1)^kc_\nu^k$，其中
$\sigma_\nu=\pm1$记录分子或分母，$c_\nu\in\mathbb Q$为斜率。
指数展开的$\varepsilon^r$项由总权重$r$的$\zeta$乘积组成；
乘$R_G$及$e^{-h\varepsilon L}$给所述界。
\end{proof}

这个证明也解释为什么需要规范：若保留$e^{-h\gamma_E\varepsilon}$，
展开就含Euler常数；这不能被省略后写成MZV定理。
去掉$x^{-A_0}$后，当$x=\mu^2$时系数属于$\mathcal Z_1$；一般$x$时是
$\mathcal Z_1[L]$，不能把任意质量单位或运动学对数自动叫作MZV常数。
有质量、非整数中心指数或内部动量注入都可能破坏这个接受域。

\begin{proposition}[有限证书的核查成本]\label{prop:masslesscertificate}
给定边表与上述二叉证书，可用不超过$O(E^2)$次集合/端点核对以及
$O(E)$个指数加减操作核实接受条件并生成至多$6h$个Gamma因子的表达树。
无需生成生成树多项式、求积或数值整数关系。
\end{proposition}
\begin{proof}
满二叉树有$2E-1$个结点。逐结点比较子图的带标号边集与粘合端点，
用朴素至多$E$长度列表即可在$O(E^2)$内验证；同时检查外动量标记。
三条递归量$(h,A,C)$每个结点只需常数次操作，每个$P$结点追加六个Gamma记录。
这些是组合操作及符号表达树的界；若把任意大整数位移完全展开成线性因子，
或列举所有$\varepsilon$阶系数，其位复杂度另计，不被这个界隐瞒。
\end{proof}

\subsection{重整化与局部RG的接受条件}
\begin{proposition}[森林数据闭合时的MZV保持]\label{prop:masslessforestclosure}
固定一个有限阶的实际MS森林计算。假设无未分离的IR极点，局部反项规则已成立，
且每个需要的子图和带局部插入的余图在去量纲后，都有上一定理的证书，
或是这类表达的有限有理组合和明确为零的无尺度积分。
则该计算的所有MS反项、有限部分、以及由它们作有限有理代数操作与尺度导数所得的
RG图贡献，系数属于$\mathcal Z_1$；保留尺度时属于$\mathcal Z_1[L]$。
这不声称任一串并联图自动满足所有森林接受条件。
\end{proposition}
\begin{proof}
有限圈阶按子图数递归，用
\[
 C_-(G)=-K\!\left(F_G+
 \sum_{\emptyset\ne\gamma\subsetneq G}C_-(\gamma)F_{G/\gamma}\right),
 \qquad R(G)=(1-K)\overline R(G),
\]
其中求和和局部插入按给定实际森林规则，$K$取负次幂。
归纳假设与接受证书使括号中的各项都属于相应Laurent环。
有限加法、乘法和$K$保持该环。若需要除以常数项为非零有理数的形式级数，
其逆也逐阶留在同一环。$\partial_L$同样保持此环；
局部性另外保证反项的局部系数不含$L$。
没有将包含IR污染的裸极点解释为UV反项，也没有由单图贡献推出完整理论的$\beta$函数。
若有限方案另添任意常数，该结论只在所添常数仍属所声明代数时保持。
\end{proof}

\subsection{实际无质量banana：全阶规则与完整物理读出}
令$B_n$为两点之间的$n\ge2$条单位无质量边，$h=n-1$。
它可以来自$n$种互异自由标量的正规序乘积源两点函数；
剥去Wick归一化和测度常数后，正是这里的有标号标量图。
不加入相同粒子的$n!$，也不把任意$n$的复合源说成同一个可重整化基本场相互作用。

\begin{theorem}[无质量banana的闭式分类与MS有限部分]\label{thm:masslessbanana}
在前述规范下，
\begin{align}
 F_{B_n}(x,\varepsilon)
 &=x^{n-2-(n-1)\varepsilon}
 \frac{\Gamma(1-\varepsilon)^n
       \Gamma(2-n+(n-1)\varepsilon)}{\Gamma(n-n\varepsilon)},
 \label{eq:masslessbananaexact}\\
 \widehat C_{B_n}(\varepsilon)
 &=R_n(\varepsilon)\exp\left[
 \sum_{k\ge2}\frac{\zeta(k)}k
 \bigl(n+(-1)^k(n-1)^k-n^k\bigr)\varepsilon^k\right],\\
 R_n(\varepsilon)&=
 \left[\prod_{j=0}^{n-2}((n-1)\varepsilon-j)
       \prod_{j=1}^{n-1}(j-n\varepsilon)\right]^{-1}.
\end{align}
该实际图在$x>0$无IR发散，所有非空真UV森林的反项插入积分在无质量DR中为零。
因此其MS反项和有限部分分别为
\begin{align}
 C_-(B_n)&=-\frac{r_n}{\varepsilon}x^{n-2},
 &r_n&=\frac{(-1)^{n-2}}{((n-1)!)^2},\\
 F^{\mathrm{MS}}_{B_n}(x)&=r_nx^{n-2}\bigl(b_n-(n-1)L\bigr),
 &b_n&=(n-1)H_{n-2}+nH_{n-1},\label{eq:masslessbananaMS}
\end{align}
其中$H_0=0$、$H_m=\sum_{j=1}^m1/j$。
特别地，去量纲后的减除点值是有理数；所有更高$\varepsilon$系数均在单$\zeta$代数中。
\end{theorem}
\begin{proof}
更一般维数的$n$线公式为
\[
 x^{(n-1)D/2-n}
 \frac{\Gamma(D/2-1)^n\Gamma(n-(n-1)D/2)}
 {\Gamma(n(D/2-1))}.
\]
$n=2$是一圈引理。把$n-1$线子图和一条边并联，
将有效指数$n-1-(n-2)D/2$与1代入卷积核，Gamma因子望远镜消去，得到归纳步。
整数维收敛域的解析指数延拓同前，故所得是实际图的亚纯表达。
把$D=4-2\varepsilon$代入，再把Gamma平移至1，给$R_n$及指数中的系数。

IR方面，非零$q$使全部$n$条线不能同时软化。
任一真子集的$r<n$个软线动量可作为独立的$r$个四维局部坐标；
其奇性至多为$\prod_{j=1}^r(k_j^2)^{-1}$，局部可积。
这同时覆盖移位后某条线变软的层；四维附近仍成立。
因此这里的零点极点是UV问题，不能套用有IR的无尺度原图。

每个真圈子图包含两端顶点和$r<n$条平行边。
把其局部反项缩成一点后，至少一条余边成为无质量蝌蚪；
局部反项只提供有限阶动量多项式。
展开这些多项式后，余图是独立无尺度张量蝌蚪积分的有限和，DR中均为零。
嵌套森林按最外层真子图同样判断；互相重叠的子图不能当作可同时收缩的森林。
所以$\overline R(B_n)=\widehat F_{B_n}$，无需把裸图一般地等同于已减除图。

$R_n$的唯一零点极点来自$j=0$，其留数为$r_n$。
Gamma指数从$\varepsilon^2$开始，而
\[
 R_n(\varepsilon)=\frac{r_n}{\varepsilon}
                 (1+b_n\varepsilon+O(\varepsilon^2)),\qquad
 e^{-h\varepsilon L}=1-h\varepsilon L+O(\varepsilon^2).
\]
二者相乘即给反项及有限部分，$b_n$由有限线性因子的对数导数求出。
这是确切全$n$公式，不依赖数值拟合。
\end{proof}

\begin{corollary}[MOM、正谱与四线显式物理函数]\label{cor:masslessphysical}
令$T_{\mu^2}^{n-2}$取$x=\mu^2$处至$n-2$阶Taylor多项式。
同一源函数满足这些局部MOM条件的代表为
\[
 F^{\mathrm{MOM}}_{B_n}(x)
 =-(n-1)r_n(1-T_{\mu^2}^{n-2})
       \bigl[x^{n-2}\log(x/\mu^2)\bigr].
\]
取$\mathcal A(s\pm i0)=F_{B_n}^{\mathrm{MS}}(-s\mp i0)$，则$s>0$时
\[
 \frac{\operatorname{Disc}\mathcal A(s)}{2\pi i}
 =\frac{s^{n-2}}{(n-1)!(n-2)!}>0.
 \label{eq:masslessspectral}
\]
同一不连续性也适用于上述MOM代表。具体四线有
\begin{align}
 F^{\mathrm{MS}}_{B_4}(x)&=\frac{x^2}{36}
                   \left(\frac{71}{6}-3\log(x/\mu^2)\right),\\
 F^{\mathrm{MOM}}_{B_4}(x)&=-\frac1{12}\left[
 x^2\log(x/\mu^2)-\mu^2(x-\mu^2)-\frac32(x-\mu^2)^2\right],\\
 \frac{\operatorname{Disc}\mathcal A_4(s)}{2\pi i}&=\frac{s^2}{12}.
\end{align}
这些是前述剥去常数的源两点函数，不是未经LSZ和通量处理的散射截面。
\end{corollary}
\begin{proof}
两种局部延拓之差是次数至多$n-2$的多项式；指定Taylor喷射唯一确定它。
MS有限部分中的$b_nx^{n-2}$全部被Taylor减除，得到首式。
对于$s>0$，$\log(-s\mp i0)=\log s\mp i\pi$，
而$(-s\mp i0)^{n-2}=(-1)^{n-2}s^{n-2}$。
代入MS式与$r_n$，符号相消，得到正谱公式。
局部多项式没有割，故MOM同样成立。
最后用$r_4=1/36$、$b_4=71/6$及
$T_{\mu^2}^{2}[x^2\log(x/\mu^2)]=\mu^2(x-\mu^2)+3(x-\mu^2)^2/2$。
\end{proof}

\subsection{同一拓扑为何从K3退化到对数}
\begin{proposition}[四线全零质量层的显式有理化]\label{prop:masslessdegeneration}
固定$s\ne0$，四线谱载体在$m_1=m_2=m_3=m_4=0$时为
\[
 y^2=a^2b^2\lambda(s,a,b),\qquad
 \lambda(s,a,b)=(s-a-b)^2-4ab.
\]
约去重复因子后，在$ab\ne0$的稠密开集上有理坐标
\[
 z=y/(ab),\quad t=\frac{s-a-b-z}{2a},\qquad
 b=\frac{st}{1+t}-at,\quad z=\frac{s}{1+t}-a(1+t).
\]
因此一般有质量的K3载体在此奇异质量层的正规化是有理的；
实际无质量四维函数的对数分类由上面的完整物理公式确定。
\end{proposition}
\begin{proof}
把$t$的定义写成$z=s-a-b-2at$，代入$z^2=(s-a-b)^2-4ab$，
约去非零$4a$即得$(1+t)b=st-at(1+t)$。
这给互逆有理参数化，$t=-1$及$a b=0$只需作为边界保留。
有理化本身不分类全部相对扩张；上一定理和推论另行给完整读出的对数形式。
所以这里没有从环境有理性直接跳到完整周期结论。
\end{proof}

这一结果与分类表K3行保持同一$D=4$、同一四线拓扑，改变的是质量层和已声明的减除输入。
全零质量必须作为M4独立一行，不能从光滑正质量族把K3标签机械外推。
反过来，也不能在正质量积分中直接丢掉质量后保留原答案。
本节独立定义并解析延拓无质量图，不假定质量极限、维数展开及减除三者无条件可交换。
在$D=2$直接全零质量会遇到新的IR问题，因而不在此物理MOM结论内。

表的M3行由\ref{thm:masslesssp}和\ref{prop:masslesscertificate}给全阶肯定证据；
森林物理出口须提交\ref{prop:masslessforestclosure}的附加输入。
M4行由\ref{thm:masslessbanana}--\ref{prop:masslessdegeneration}实际闭合MS、MOM和正谱。
这回答了为什么某些周期能用MZV，以及怎样在积分前给充分判断；
它不提供所有图的必要判据，也不识别每个正规化系数原生的完整最小动机。

\begin{thebibliography}{9}
\bibitem{PanzerMZV} E. Panzer, \emph{Feynman integrals and hyperlogarithms},
arXiv:1506.07243; author thesis, equation (1.2.1), \S5.2.1.
\url{https://people.maths.ox.ac.uk/panzer/paper/phd.pdf}.
One-loop convolution and one-scale insertion rules are classical input; the present
classification certificate and the displayed normalization are derived explicitly.
\bibitem{DLMFGamma} NIST Digital Library of Mathematical Functions, \S5.7, equation 5.7.3,
\url{https://dlmf.nist.gov/5.7.E3}.
\end{thebibliography}
\end{document}
